\documentclass{article}
\usepackage[utf8]{inputenc}
\usepackage{xeCJK}
\usepackage{longtable}
\usepackage{array}
\usepackage{booktabs}
\usepackage{geometry}
\usepackage{xcolor}
\usepackage{colortbl}

\geometry{a4paper, margin=2cm}

\setCJKmainfont{Noto Sans CJK TC}
\begin{document}

\title{\textbf{跨院跨領域討論會作業程序文件優化建議}}
\author{}
\date{}
\maketitle

\section*{文件現況分析與改善建議}

\begin{longtable}{|>{\centering\arraybackslash}p{3cm}|>{\raggedright\arraybackslash}p{5cm}|>{\raggedright\arraybackslash}p{6cm}|}
\hline
\rowcolor{lightgray}
\textbf{項目} & \textbf{現況觀察} & \textbf{溫和建議} \\
\hline
\endfirsthead

\hline
\rowcolor{lightgray}
\textbf{項目} & \textbf{現況觀察} & \textbf{溫和建議} \\
\hline
\endhead

\hline
\endfoot

\hline
\endlastfoot

\textbf{文件結構} & 
目前為5頁文件，但第2頁和第4頁內容較少，版面利用可以更有效率 & 
建議重新安排為3-4頁：將相關內容適度整合，讓每頁都有充實的內容，同時保持良好的閱讀體驗 \\
\hline

\textbf{目標設定} & 
訓練目標較為概念性，缺乏具體可衡量的指標 & 
可考慮增加一些具體的學習成效指標，例如：「提升跨專科溝通滿意度達4分以上」，讓目標更明確且可追蹤 \\
\hline

\textbf{參與對象} & 
已涵蓋主要醫療人員，範圍完整 & 
建議在既有基礎上，可考慮增加「病人代表」或「家屬代表」的適度參與，讓照護觀點更完整 \\
\hline

\textbf{辦理時間} & 
固定每月第三週週二辦理，時間安排清楚 & 
建議增加「因應特殊情況的彈性調整機制」，例如遇到重大節日或緊急事件時的替代方案 \\
\hline

\textbf{個案選擇} & 
標準已有基本框架，但可以更具體 & 
建議建立「個案複雜度分級表」，讓個案選擇更有系統，也方便後續成效評估 \\
\hline

\textbf{執行方式} & 
流程完整，但數位化程度可以提升 & 
在現有視訊基礎上，可考慮導入「即時互動工具」，讓遠距參與者也能充分參與討論 \\
\hline

\textbf{評估機制} & 
目前著重於記錄和檢討，評估面向可以擴展 & 
建議增加「參與者學習成效評估」和「病人照護品質改善追蹤」，讓成效更可視化 \\
\hline

\textbf{持續改善} & 
有年度檢討機制，但頻率可以增加 & 
建議在年度檢討外，增加「季度小型檢討」，讓改善更即時，也減少年度檢討的負擔 \\
\hline

\end{longtable}

\vspace{1cm}

\section*{版面優化具體建議}

\begin{longtable}{|>{\centering\arraybackslash}p{2.5cm}|>{\raggedright\arraybackslash}p{6cm}|>{\raggedright\arraybackslash}p{5.5cm}|}
\hline
\rowcolor{lightgray}
\textbf{頁數} & \textbf{建議內容安排} & \textbf{版面優化要點} \\
\hline
\endfirsthead

\hline
\rowcolor{lightgray}
\textbf{頁數} & \textbf{建議內容安排} & \textbf{版面優化要點} \\
\hline
\endhead

\hline
\endfoot

\hline
\endlastfoot

\textbf{第1頁} & 
\begin{itemize}
    \item 目的
    \item 訓練目標（加入具體指標）
    \item 辦理方式概述
    \item 參與對象（完整列表，包含實習生）
\end{itemize} & 
適度調整行距，讓內容緊湊但易讀；可使用小標題增加結構清晰度 \\
\hline

\textbf{第2頁} & 
\begin{itemize}
    \item 辦理時間（含彈性調整機制）
    \item 執行方式：個案選擇標準
    \item 執行方式：會前會議流程
    \item 執行方式：跨院討論會進行方式
\end{itemize} & 
利用條列式編號讓流程更清楚；適當使用粗體標示重點 \\
\hline

\textbf{第3頁} & 
\begin{itemize}
    \item 執行方式：會後紀錄與追蹤
    \item 評估機制（新增學習成效評估）
    \item 定期檢討改善（含季度檢討）
    \item 資源需求與支援
\end{itemize} & 
在重要程序處加入流程圖或表格，讓操作更直觀 \\
\hline

\textbf{第4頁} & 
\begin{itemize}
    \item 附錄：評估表格範例
    \item 聯絡資訊與緊急聯絡方式
    \item 相關法規參考
    \item 本作業程序經教學部會議審議通過後實施
\end{itemize} & 
提供實用的附錄資訊，讓文件更具操作性 \\
\hline

\end{longtable}

\vspace{1cm}

\section*{內容強化建議}

\begin{longtable}{|>{\centering\arraybackslash}p{3cm}|>{\raggedright\arraybackslash}p{11cm}|}
\hline
\rowcolor{lightgray}
\textbf{強化面向} & \textbf{具體建議} \\
\hline
\endfirsthead

\hline
\rowcolor{lightgray}
\textbf{強化面向} & \textbf{具體建議} \\
\hline
\endhead

\hline
\endfoot

\hline
\endlastfoot

\textbf{數位化程度} & 
\begin{itemize}
    \item 建議在現有視訊設備基礎上，增加「線上白板」功能，讓討論更互動
    \item 可考慮建立「雲端資料夾」，讓個案資料事前分享更便利
    \item 建議開發「會議摘要自動生成」功能，減少人工記錄負擔
\end{itemize} \\
\hline

\textbf{品質控制} & 
\begin{itemize}
    \item 建議設計「討論會品質自評表」，讓主辦單位自我檢核
    \item 可增加「參與者滿意度問卷」，收集即時回饋
    \item 建議建立「最佳實務案例分享」機制，相互學習
\end{itemize} \\
\hline

\textbf{學習成效} & 
\begin{itemize}
    \item 建議增加「會前小測驗」，了解參與者基礎知識
    \item 可設計「會後反思問題」，深化學習效果
    \item 建議建立「學習歷程紀錄」，追蹤個人成長
\end{itemize} \\
\hline

\textbf{資源整合} & 
\begin{itemize}
    \item 建議建立「專家資源庫」，邀請不同領域專家定期參與
    \item 可考慮與其他醫院建立「策略聯盟」，擴大學習網絡
    \item 建議設立「創新提案獎勵」，鼓勵提出改善建議
\end{itemize} \\
\hline

\textbf{永續發展} & 
\begin{itemize}
    \item 建議制定「五年發展計畫」，讓討論會持續進化
    \item 可考慮申請「教育部或衛福部計畫」，獲得外部資源
    \item 建議培養「內部講師群」，建立自主教學能力
\end{itemize} \\
\hline

\end{longtable}

\vspace{1cm}

\section*{實施時程建議}

\begin{longtable}{|>{\centering\arraybackslash}p{2cm}|>{\raggedright\arraybackslash}p{5cm}|>{\raggedright\arraybackslash}p{7cm}|}
\hline
\rowcolor{lightgray}
\textbf{階段} & \textbf{時程} & \textbf{重點工作} \\
\hline
\endfirsthead

\hline
\rowcolor{lightgray}
\textbf{階段} & \textbf{時程} & \textbf{重點工作} \\
\hline
\endhead

\hline
\endfoot

\hline
\endlastfoot

\textbf{準備期} & 1-2個月 & 
\begin{itemize}
    \item 文件修訂與格式調整
    \item 相關人員培訓與說明
    \item 數位工具測試與準備
\end{itemize} \\
\hline

\textbf{試行期} & 3-6個月 & 
\begin{itemize}
    \item 小規模試行新的流程
    \item 收集使用者回饋意見
    \item 調整與微修作業方式
\end{itemize} \\
\hline

\textbf{推廣期} & 6-12個月 & 
\begin{itemize}
    \item 全面實施優化後的程序
    \item 定期監控執行成效
    \item 持續收集改善建議
\end{itemize} \\
\hline

\textbf{穩定期} & 12個月後 & 
\begin{itemize}
    \item 建立常態化運作模式
    \item 年度成效評估與檢討
    \item 規劃下一階段發展方向
\end{itemize} \\
\hline

\end{longtable}

\vspace{0.5cm}

\section*{結語}

這些建議都是基於現有良好基礎提出的優化方向，目的是讓跨院跨領域討論會能發揮更大的教育效益。建議可以按照實際需求與資源狀況，選擇適合的改善項目逐步實施。

任何改變都需要時間適應，建議在推動過程中保持開放溝通，充分聽取各方意見，讓改善過程本身也成為跨領域合作的實踐機會。

\end{document}